%! AP Statistics — Unit 1, Lesson 1.6 — Guided Notes ANSWER KEY
\documentclass[10pt]{article}
\usepackage{apstats-article}
\usepackage{apstats-key}

%=============================================================================
\begin{document}

\pageheader{Unit 1, Lesson 1.6}{Guided Notes --- Answer Key}
\begin{tabularx}{\linewidth}{X X X}
Name:\;\ans{Key} & Date:\; & Period:\;
\end{tabularx}

\vspace{0.12in}

\begin{objectivebox}
\small By the end of this lesson, I will be able to\dots\\[4pt]
\ansline{identify the shape of a distribution and describe it using SOCS.}\\[1pt]
\ansline{write a complete SOCS description with context in every sentence.}\\[1pt]
\ansline{explain the role of outliers and their effect on center and spread.}
\end{objectivebox}

\vspace{0.1in}

\begin{vocabbox}
\noindent\textbf{\textcolor{navy}{SOCS:}}\quad\ans{The AP framework for describing
quantitative distributions: Shape, Outliers, Center, Spread.}\\[4pt]

\noindent\textbf{\textcolor{navy}{Shape:}}\quad\ans{The overall pattern of a
distribution: symmetric, skewed right/left, unimodal, bimodal, uniform.}\\[4pt]

\noindent\textbf{\textcolor{navy}{Skewed right:}}\quad\ans{Long tail toward
higher (right) values; bulk of data clusters at lower values. Mean $>$ median.}\\[4pt]

\noindent\textbf{\textcolor{navy}{Skewed left:}}\quad\ans{Long tail toward
lower (left) values; bulk of data clusters at higher values. Mean $<$ median.}\\[4pt]

\noindent\textbf{\textcolor{navy}{Symmetric:}}\quad\ans{Roughly mirror-image
about the center; mean $\approx$ median.}\\[4pt]

\noindent\textbf{\textcolor{navy}{Unimodal / Bimodal / Uniform:}}\quad
\ans{One peak / Two peaks / All bars roughly equal height.}\\[4pt]

\noindent\textbf{\textcolor{navy}{Outlier:}}\quad\ans{An individual value that
falls far from the overall pattern of the distribution; investigate, do not
automatically remove.}
\end{vocabbox}

\vspace{0.1in}

\begin{hookbox}
\small\textbf{Two Histograms}

\begin{enumerate}[leftmargin=1.5em, itemsep=6pt]
  \item \ans{Skewed right / right-skewed / ``bunched at the low end.''
        Household incomes have a few extremely high earners pulling the tail
        to the right.}

  \item \ans{Symmetric / bell-shaped / unimodal. Heights cluster near the center
        (around 64 inches) with similar drop-off on both sides.}

  \item \ans{The shape matters because it affects which summary statistics are
        most appropriate. For a skewed distribution, the median is a better
        center than the mean; for symmetric data, both work. Shape also suggests
        what real-world process generated the data.}
\end{enumerate}
\end{hookbox}

\vspace{0.1in}

\begin{notesbox}{The SOCS Framework}
\small

\textbf{\textcolor{navy}{S --- Shape}} \quad
Overall pattern: \ans{symmetric}, skewed right, skewed left.
Number of peaks: \ans{unimodal} or \ans{bimodal}.
Special case: \ans{uniform}.

\vspace{4pt}
\textbf{\textcolor{navy}{O --- Outliers}} \quad
Values that fall \ans{far from the overall pattern}.
Note their approximate \ans{value} and how far they are from the bulk.
Always \ans{investigate} outliers.

\vspace{4pt}
\textbf{\textcolor{navy}{C --- Center}} \quad
Informal estimate of a \ans{typical} value.
For now: use the \ans{median} or approximate balance point.
Lesson 1.7: \ans{mean} and \ans{median}.

\vspace{4pt}
\textbf{\textcolor{navy}{S --- Spread}} \quad
How \ans{wide} or \ans{narrow} the distribution is.
For now: describe the \ans{range} (min to max).

\vspace{4pt}
\textbf{AP Context Rule:} Every sentence must name the \ans{variable} and the \ans{population}.

\end{notesbox}

\newpage

\begin{notesbox}{Shape Vocabulary}
\small

\renewcommand{\arraystretch}{2.0}
\begin{tabularx}{\linewidth}{@{} >{\bfseries\color{navy}}p{0.22\linewidth}
                                   p{0.38\linewidth}
                                   X @{}}
\rowcolor{sky}
\textbf{Shape} & \textbf{What it looks like} & \textbf{Real-world example} \\
\midrule
Symmetric      & Mirror image on each side & \ans{Heights, IQ scores, many measurement errors} \\
Skewed right   & Long tail toward \ans{higher} values & \ans{Incomes, house prices, wait times} \\
Skewed left    & Long tail toward \ans{lower} values & \ans{Age at retirement, test scores on easy tests} \\
Unimodal       & \ans{One} clear peak & \ans{Heights, weights} \\
Bimodal        & \ans{Two} clear peaks & \ans{Ages at a concert with two audience groups} \\
Uniform        & All bars roughly \ans{equal} height & \ans{Rolling a fair die, random number tables} \\
\end{tabularx}

\vspace{6pt}
\textbf{Common error:} The skew direction = direction the \ans{tail} points.

\end{notesbox}

\vspace{0.1in}

\begin{notesbox}{Worked Example --- SAT Math Scores}
\small

Model SOCS sentence labels: S, O, C, S (in order of the sample response).

\vspace{6pt}
Strong AP response elements:
\begin{enumerate}[itemsep=3pt]
  \item Context in the \ans{first} sentence (and throughout).
  \item All \ans{four} SOCS components present.
  \item Outliers \ans{hedged} (``appear to be'').
  \item Center and spread include \ans{units} (``points'').
\end{enumerate}

\end{notesbox}

\vspace{0.1in}

\begin{practicebox}
\small\textbf{Daily text messages --- 40 students}

\begin{enumerate}[leftmargin=1.5em, itemsep=6pt]

  \item \ans{Roughly bimodal or skewed right. Main cluster at 20--60 messages;
        a secondary cluster at 100--140 suggests a second group. The long tail
        to the right supports a skewed right interpretation. Either answer is
        acceptable with justification.}

  \item \ans{The student at 180 messages is a clear outlier --- far above the
        next highest values (100--140 range). The 100--140 cluster may also
        warrant investigation depending on how it is defined.}

  \item \ans{Center: approximately 45 messages per day (near the median of the
        main cluster). Spread: 5 to 180 messages per day, a range of 175 messages.}

  \item \ans{The distribution of daily text messages for these 40 students is
        roughly bimodal and skewed right, with a large cluster between 20 and
        60 messages and a smaller cluster between 100 and 140 messages. [S]
        One student sent 180 messages per day, a potential outlier that falls
        well above the rest of the group. [O] The center of the distribution
        is approximately 45 messages per day, representing a typical student
        in this group. [C] The spread is large, ranging from 5 to 180 messages
        per day, a range of 175 messages. [S]}

\end{enumerate}
\end{practicebox}

\vspace{0.1in}

\begin{spiralbox}
\small
\textbf{Big Idea:} \ans{Saying ``skewed right'' alone is insufficient for AP
credit. Graders also require: outliers addressed (even if none), center estimated
with units and context, spread described (at minimum as a range), and context in
every sentence. Missing any one component costs points on the AP exam.}
\end{spiralbox}

\end{document}
